\section{Implementation}\label{sec:implementation}

The proposed forecasting framework is implemented as a Python-based
software package consisting of \textbf{eight core modules}, an
integrated research notebook, and a \textbf{Flask web application}. The
research notebook executes the complete experimental pipeline from data
preparation to evaluation, while the Flask application provides an
interactive interface that uses the same underlying modules. This shared
implementation ensures that both the experimental results and the
deployed application produce consistent outputs.

The entire system is developed using \textbf{Python 3.12}. The machine
learning component is implemented using \textbf{Keras 3.15} with a
\textbf{PyTorch backend}, eliminating the need for a TensorFlow
installation. Data preprocessing, feature scaling, TF-IDF-based evidence
retrieval, and the classical machine learning algorithms are implemented
using \textbf{scikit-learn 1.9}. The language-model components
communicate with \textbf{Gemini 2.5 Flash} through the
\textbf{google-genai} client library. Data visualization is performed
using \textbf{Matplotlib}, while the web interface is developed using
\textbf{Flask 3.1}.

To ensure that the reported experimental results can be reproduced, the
software versions of all libraries are fixed throughout the
implementation. Using the same package versions guarantees that the
models, evaluation metrics, and reported performance remain consistent
across repeated executions of the framework.

\begin{table}[H]
\centering
\footnotesize
\begin{tabularx}{\textwidth}{@{}p{2.8cm}Y>{\raggedright\arraybackslash}p{3.6cm}@{}}
\toprule
\textbf{Module} & \textbf{Responsibility} & \textbf{Key output} \\
\midrule
\texttt{config.py} & Every experimental setting in one place, imported by both the notebook and the web app & -- \\
\texttt{data\_loader.py} & Price and VIX loading, the 22 backward-looking features, the 5-day label & feature frames \\
\texttt{ml\_model.py} & Walk-forward LSTM training, embargo, threshold selection, serving & \texttt{ml\_predictions.csv} \\
\texttt{evidence.py} & Builds the 252-item dated corpus; TF-IDF retrieval with the look-ahead guard & \texttt{evidence\_corpus.json} \\
\texttt{llm.py} & Prompt construction, Gemini calls, response parsing, repeated-run self-consistency & \texttt{llm\_responses.csv} \\
\texttt{hybrid.py} & Calibrated logit pooling, regime weight derivation, risk-profile advice & \texttt{hybrid\_weights.json} \\
\texttt{evaluation.py} & Three hallucination checks, McNemar, date-blocked bootstrap, rubric scoring & \texttt{results.json} \\
\texttt{assistant.py} & Question routing and grounded free-text answering for the web assistant & -- \\
\texttt{app.py} & Flask routes serving all four arms side by side, plus a live stability tab & web interface \\
\bottomrule
\end{tabularx}
\caption{The nine implementation modules and what each produces.}
\label{tab:implementation-modules}
\end{table}

A key design principle of the implementation is that both the research
notebook and the Flask web application use the \textbf{same software
modules and configuration settings}. No separate implementation is
created for demonstration purposes. As a result, every prediction,
threshold, and evaluation result displayed in the web interface is
generated from the same codebase used during experimentation. This
unified design eliminates inconsistencies between the deployed
application and the reported research findings. In earlier versions of
the system, separate threshold settings were used for the demonstration
interface, which occasionally produced predictions that differed from
those reported in the experimental results. This issue was removed by
adopting a single shared implementation.

\textbf{Reproducibility}

Reproducibility is an important objective of the proposed framework. To
ensure consistent experimental results, every response generated by
\textbf{Gemini 2.5 Flash} is cached locally using a unique identifier
created from the prompt content and execution index. This caching
mechanism allows all \textbf{1,000 responses} used in the hallucination
analysis and all \textbf{150 responses} used in the stability
experiments to be reproduced without issuing additional API requests.

Randomness is also carefully controlled throughout the implementation.
Fixed random seeds are applied to \textbf{Python}, \textbf{NumPy}, and
\textbf{Keras}, while the fold index is incorporated into the seed value
so that individual folds remain reproducible but are not identical.
Executing the research notebook with its default configuration
automatically regenerates all experimental tables and figures using the
stored prediction files, allowing the complete evaluation to be
reproduced in less than one minute.

To maintain data integrity and security, the Gemini API key is stored
separately in an environment configuration file rather than within the
source code. If the required API credentials are unavailable, the
execution process terminates instead of generating placeholder outputs,
ensuring that every reported result is based on valid experimental data.

